\begin{abstract}
摘要部分用于简要介绍研究工作的主要内容、研究方法、主要结论和创新点。
摘要应独立成文，使读者能够快速了解报告的核心内容。

本示例模板展示博士后研究报告的摘要格式规范。摘要一般应包含以下几个方面：

（1）研究背景与意义：简要说明研究工作的背景、目的和意义。

（2）研究内容与方法：概述研究的主要内容、采用的研究方法和技术路线。

（3）主要结论与成果：总结研究的主要发现、结论和创新成果。

（4）应用价值与展望：说明研究成果的应用价值及未来研究方向。

摘要篇幅建议控制在500-1000字左右，内容应简明扼要、重点突出。
避免使用过于专业的术语，使摘要具有较强的可读性。

\keywords{关键词1，关键词2，关键词3，关键词4，关键词5}

\begin{englishabstract}
    The abstract section is used to briefly introduce the main content, 
    research methods, major conclusions, and innovations of the research work. 
    The abstract should be self-contained, enabling readers to quickly 
    understand the core content of the report.
    
    This template demonstrates the format requirements for the abstract 
    of a postdoctoral research report. The abstract generally should 
    include the following aspects:
    
    (1) Research background and significance: briefly explain the background, 
    purpose, and significance of the research work.
    
    (2) Research content and methods: outline the main research content, 
    research methods adopted, and technical approach.
    
    (3) Major conclusions and achievements: summarize the main findings, 
    conclusions, and innovative results of the research.
    
    (4) Application value and prospects: explain the application value 
    of research results and future research directions.
    
    The abstract length is recommended to be around 500-1000 words, 
    with concise and focused content. Avoid using overly specialized 
    terminology to ensure readability.
    
    \englishkeywords{Keyword 1, Keyword 2, Keyword 3, Keyword 4, Keyword 5}
\end{englishabstract}
\end{abstract}